% backmatter/supplement/s-ch1.tex -- Mandell's ``Changes to Text of Chapter 1'' (June 1994): the cover note
% (PDF 667, signed ``Gordon Mandell / June 1994''), the changes to pages 14, 16, 18, 19 and 30-31 (PDF 668-672),
% the symbols to add (PDF 673) and the new Figure 2 (PDF 674), shown with the 1973 Figure 2 (PDF 44); then the
% ``Correction to Original Page 40 of Chapter 1'' (PDF 675). All typewritten.
% The pressure term is set as these pages write it, with the arrow over the whole term,
% \overrightarrow{(P_e - P_a)A_e} (his note of 1 June 1994, supp:vector, later asks for (P_e - P_a)\vec{A}_e).
% The typed arrows over A_e are set \vec{A}_e, the form the chapters use for this symbol. The equation numbers,
% printed at the left, are the documents' own: \tag without \label (the chapter owns the labels). The two
% caption displays are gather* so that the long arrow keeps the full display skip below the short line above.
\chapter{Changes to Text of Chapter 1}\label{supp:ch1}

The equation for rocket thrust that originally appeared in \emph{Topics in Advanced Model Rocketry} was
restricted to the case where the rocket's exhaust expands through the divergent section of the nozzle until
it crosses the nozzle exit plane at a pressure equal to the ambient atmospheric pressure. This condition is
called ``optimal expansion'' or ``correct expansion'' and many model rocket motors closely approach it. Some
readers were, however, dissatisfied with the restriction so I wrote these changes to the text of
Chapter~\ref{ch1} to include cases where the exhaust pressure is greater or less than the ambient pressure.

\bigskip
\noindent Gordon Mandell\\
June 1994

\bigskip
\begin{center}
[Change caption of Figure~\ref{ch1:fig:2} on page 14 to read as follows]:
\end{center}

\noindent\textbf{Figure 2:} Origin of Rocket Thrust. In a time interval $\Delta t$ the rocket expels the
quantity of mass $\Delta m_e$ at the exhaust velocity $\vec{c}$ in the $(-y)$ direction. ``In the limit'' as
$\Delta t$ becomes $dt$ and $\Delta m_e$ becomes $dm_e$ the rate of mass expulsion---or ``mass flow
rate''---is $dm_e/dt$ (also written $\mdot$). The product of the mass flow rate and the exhaust velocity
produces a component of force $-\vec{c}(dm_e/dt)$ acting in the $(+y)$ direction. The exhaust gases leave
the rocket at the exhaust pressure $P_e$ when they cross the plane at the aft end of the rocket motor
nozzle, called the \emph{nozzle exit plane}. The difference between $P_e$ and the pressure of the atmosphere
outside the exhaust stream (also called the \emph{ambient} pressure) $P_a$ acts over the area of the nozzle
exit plane $\vec{A}_e$ to produce another component of force $\overrightarrow{(P_e - P_a)A_e}$ which also
acts in the $(+y)$ direction when $P_e$ is greater than $P_a$, but which acts in the $(-y)$ direction when
$P_e$ is less than $P_a$. The sum of the two force components is the rocket's \emph{thrust} $\vec{F}$:
\begin{gather*}
\vec{F} = -\vec{c}(dm_e/dt) + \overrightarrow{(P_e - P_a)A_e}
\end{gather*}
\emph{NOTE}: In order to make the illustrations easier to see, the exhaust stream has been drawn as if the
nozzle exit diameter were equal to the body tube diameter. The nozzle exit of most model rocket engines is
actually much smaller than the body tube diameter.

\begin{center}
[Insert after the 3rd line on page 16]:
\end{center}

As the exhaust gases leave the nozzle they cross the plane at its aft end, called the \emph{nozzle exit
plane}, at a pressure called the \emph{exhaust pressure} and enter the external environment where the
pressure everywhere outside the exhaust stream is the pressure of the surrounding atmosphere, also called the
\emph{ambient} pressure.

\begin{center}
[Change the 4th and subsequent lines on page 16 to read as follows]:
\end{center}

Now the thrust of a rocket engine is related to the time derivative of the rocket's mass and to the velocity
and pressure at which the engine is ejecting exhaust gases by the fact that \emph{the thrust of a rocket
motor is equal to the rate at which it is transferring momentum to the rocket in which it is mounted}. Those
of you who have had some exposure to elementary physics will be familiar with the concept of momentum; some
of you, moreover, may already know that the rate of momentum transfer resulting from the expulsion of mass
from a body depends upon the \emph{rate} at which mass is being expelled and the \emph{velocity} with which
it is being ejected, and that the rate of momentum transfer resulting from the application of a force to a
body is just equal to the force. In Chapter~\ref{ch4} it will be shown by more detailed momentum
considerations that the relationship involved is
\begin{equation*}\tag{1}
\vec{F} = -\vec{c}(dm_e/dt) + \overrightarrow{(P_e - P_a)A_e}
\end{equation*}
\begin{itemize}[nosep,leftmargin=4em,labelwidth=3.5em,labelsep=0.5em,align=left]
  \item[where] $\vec{F}$ = thrust produced by the rocket motor
  \item[] $\vec{c}$ = velocity of the exhaust gases \emph{as measured by an observer moving with the
    rocket}.
  \item[] $dm_e/dt$ = rate at which rocket motor is expelling mass
  \item[] $P_e$ = exhaust pressure
  \item[] $P_a$ = ambient pressure
  \item[] $\vec{A}_e$ = area of the rocket nozzle's exit plane
\end{itemize}
Students of physics will note that the total thrust, the nozzle exit plane area and the pressure component
of thrust, and the exhaust velocity are \emph{vector} quantities; that is, both their magnitudes and the
directions in which they act must be specified to describe them completely. Mass and its derivative with
respect to time, on the other hand, are \emph{scalar} quantities which are completely described by their
magnitudes and algebraic signs.

In a properly designed and constructed model rocket whose engine is not suffering from any malfunction the
exhaust stream
\begin{center}
[End of page 16 and beginning of page 18 of original text]
\end{center}
is directed dead astern. In such a case we can replace the full vector notation for forces and velocities
with a simpler, fore-and-aft sign convention in which vector quantities directed forward along the vehicle
longitudinal axis are considered algebraically \emph{positive}, while those directed rearward are considered
algebraically \emph{negative}. According to this convention,
\begin{gather*}
c = -\vec{c}\\
F = \vec{F}
\end{gather*}
and
\begin{equation*}
A_e = \vec{A}_e \qquad\text{so that}\qquad (P_e - P_a)A_e = \overrightarrow{(P_e - P_a)A_e}
\end{equation*}
Then we can write
\begin{equation*}\tag{2}
F = c(dm_e/dt) + (P_e - P_a)A_e
\end{equation*}
Professional rocket engineers refer to the derivative $dm_e/dt$ as the \emph{mass flow rate} and sometimes
write it as $\mdot$, the superscript dot denoting differentiation with respect to time (an alternative
calculus notation for the time derivative which is sometimes used to save writing). Any rocket whose exhaust
is directed dead astern will produce a thrust directed forward along the rocket's centerline, but the
pressure component can appreciably alter the magnitude of the total thrust. Most practical rocket engines
generate combustion chamber temperatures and pressures high enough to produce supersonic exhaust gas
velocities, and they incorporate convergent/divergent nozzles which efficiently convert the
combustion-chamber temperatures and pressures of their gaseous combustion products into such velocities. A
rocket engine of this type will attain the greatest possible thrust for a given combustion chamber
temperature and pressure and a given nozzle \emph{throat area} (smallest cross-sectional area within the
nozzle) when the nozzle is designed to make $P_e = P_a$ so that the pressure term of
equation~\eqref{ch1:eq:2} vanishes. This condition is known as ``optimum expansion'' or ``correct
expansion''. If the nozzle ``overexpands'' the exhaust, so that $P_e$ is less than $P_a$, the pressure term
opposes the mass expulsion term and reduces the thrust. If the nozzle ``underexpands'' the exhaust, so that
$P_e$ is greater than $P_a$, the pressure term adds to the mass expulsion term. For this reason,
$(P_e - P_a)$ is sometimes called the ``overpressure'' of an underexpanded exhaust. An underexpanding
nozzle, however, does not produce as high an exhaust velocity as one designed for optimum expansion and a
detailed examination of the gasdynamics involved shows that an underexpanding supersonic nozzle always loses
more thrust in reduction of the exhaust velocity than it gains in overpressure.

\begin{center}
[End of text on page 18]

\medskip
[Change page 19 of original text to the following:]
\end{center}

In theory, then, if you were given the exhaust velocity, exhaust pressure, and mass flow rate of any given
rocket motor, and the ambient pressure in which the motor was operating, as functions of time you could
compute the thrust of that motor as a function of time. It is also theoretically possible to calculate the
exhaust velocity and pressure and the mass flow rate from the combustion characteristics and thermodynamic
properties of the propellant and the size and shape of the combustion chamber and the nozzle. In practice,
however, rocket engineers generally compute the thrust directly from the propellant, combustion chamber, and
nozzle characteristics and then obtain the exhaust velocity and pressure and the mass flow rate as
byproducts of this calculation. The thrusts of model rocket engines as functions of time are always
supplied by the manufacturers in graphical form: the so-called \emph{thrust-time curves}, which are derived
from averaging a large number of traces recorded during static firings of engines of the same type. The
thrust-time curves of individual model rocket motors may differ somewhat from the average for their type.
This variation is not too severe, however, in motor types which have received Safety Certification or
Contest Certification from the Standards and Testing Committee of the National Association of Rocketry.

In any given engine, one generally finds that variations in thrust with time are due more to variations in
the mass flow rate than to variations in the exhaust velocity and pressure. In fact, it is almost always
accurate enough for model rocketry purposes to assume that the exhaust velocity is constant and the exhaust
is optimally expanded during the entire burning period. The mass flow rate can then be easily computed by
writing equation~\eqref{ch1:eq:2} (now simplified by the absence of the pressure term, since optimum
expansion assumes $P_e = P_a$) as
\begin{equation*}\tag{2a}
\mdot = (dm_e)/(dt) = F/c
\end{equation*}
if you are careful with units. Suppose, for

\begin{center}
[End of text on page 19]\\
{[Pages 20 through 29 are unchanged]}

\medskip
[On pages 30 and 31 of the original text, change text starting\\
with the sentence just after the question, ``How powerful is it?''\\
to the following:]
\end{center}

\noindent The answer to this question lies in a quantity called the \emph{specific impulse} of the
propellant; that is, the total impulse obtainable from the combustion of a unit weight of propellant.
Specific impulse is thus a measure of a propellant's efficiency. The definition of specific impulse is
\begin{equation*}\tag{5}
\Isp = I_t/w_f = I_t/(m_f g)
\end{equation*}
where $m_f$ is the \emph{mass} of propellant consumed in providing the total impulse $I_t$, $w_f$ is the
\emph{weight} of propellant consumed in providing that total impulse, and $g$ is the acceleration of the
Earth's gravity field at mean sea level (9.806 meters/sec.$^{2}$). According to this definition $\Isp$ has
units of newton-seconds per newton---or simply \emph{seconds}. In fact, it has units of seconds and is
numerically identical in all common systems of units. The specific impulse is related to the average
exhaust velocity over the burning period by the equation
\begin{equation*}\tag{6}
c + (P_e - P_a)A_e/\mdot = g\Isp
\end{equation*}
where the quantity $c + A_e[(P_e - P_a)/\mdot]$ is defined as the ``effective'' or ``equivalent'' exhaust
velocity $c_{\mathrm{eff}}$:
\begin{align*}
c_{\mathrm{eff}} &= c + (P_e - P_a)A_e/\mdot \tag{7}\\
&= g\Isp
\end{align*}
The specific impulse measured when a given type of propellant is burned in an actual rocket engine thus
depends not only on the propellant itself, but also to some degree on the ambient pressure and the nozzle
design. A practical rocket engine with a supersonic exhaust will achieve the highest value of $\Isp$ for a
given ambient pressure when the exhaust is optimally expanded; that is, when $P_e = P_a$ so that the pressure
term vanishes and $c_{\mathrm{eff}} = c$. The highest $\Isp$ a given type of propellant can produce under any
conditions is called the ``ideal'' or ``ideal vacuum'' specific impulse. It is the specific impulse which
would be measured if the exhaust were optimally expanded to a perfect vacuum ($P_e = P_a = 0$) through a
frictionless nozzle which converts all the energy of combustion that manifests as the temperature and
pressure of the gases in the combustion chamber into velocity. The associated exhaust velocity is called the
``ideal'' or ``ideal vacuum'' exhaust velocity. The ``ideal'' or ``ideal vacuum'' specific impulse and
exhaust velocity are theoretical maximum limits. They cannot ever quite be attained in practice, even in
space, because they would require an infinitely long frictionless nozzle with an infinite exit plane area.

Model rocket engines using pressed blackpowder propellant grains exhibit values of specific impulse in the
range of 48 to 103 seconds when operated at standard sea level ambient pressure. More recently, several
types of engines using high-performance, ammonium perchlorate/polyurethane propellants have been introduced
into model rocketry. These motors deliver specific impulses in the neighborhood of 175 seconds at sea level
ambient pressure. There exist solid propellants which deliver specific impulses in excess of 240 seconds
when burned at high combustion chamber pressures, but problems of temperature, pressure, ignition
technology, and cost have so far prevented their use in model rocketry.

The specific impulse delivered by any given model rocket motor is a quantity that is measured and
published, either by the motor manufacturer or by the NAR Standards and Testing Committee. It is therefore a
simple matter to compute the average effective exhaust velocity using equation~\eqref{ch1:eq:7}, assume
optimum expansion so $c_{\mathrm{eff}} = c$, and to use the value of $c$ thus obtained to compute the mass
flow rate as a function of time from equation~\eqref{ch1:eq:n2a} and a knowledge of the engine's thrust-time
curve. As stated previously, it is almost always sufficiently accurate for model rocket performance
calculations to assume that the exhaust velocity is constant at its average value and the exhaust is
optimally expanded during the entire burning period.

\begin{center}
[End of Section~\ref{ch1:sec:2.1}]

\medskip
[Section~\ref{ch1:sec:2.2} and remainder of Chapter~\ref{ch1} are unchanged]
\end{center}

\bigskip
\noindent\begin{minipage}{\linewidth}
The table of symbols for Chapter~\ref{ch1} will need to have the following symbols added:
\begin{center}
\begin{tabular}{@{}l p{0.72\linewidth}@{}}
\toprule
$\vec{A}_e$ & nozzle exit plane area written as a vector whose direction is forward along the vehicle
  centerline \\ \addlinespace
$A_e$ & magnitude of nozzle exit plane area \\ \addlinespace
$P_a$ & ambient pressure \\ \addlinespace
$P_e$ & pressure of exhaust gas at nozzle exit plane \\ \addlinespace
$c_{\mathrm{eff}}$ & magnitude of effective exhaust velocity (\emph{also} called equivalent exhaust
  velocity), whose direction is assumed to be rearward along the vehicle centerline \\
\bottomrule
\end{tabular}
\end{center}
\end{minipage}

\clearpage
\begin{center}
  \includegraphics{figures/v2/supplement/ch1-fig02-1994.pdf}
\end{center}

\noindent\textbf{Figure 2:} Origin of Rocket Thrust. In a time interval $\Delta t$ the rocket expels the
quantity of mass $\Delta m_e$ at the exhaust velocity $\vec{c}$ in the $(-y)$ direction. ``In the limit'' as
$\Delta t$ becomes $dt$ and $\Delta m_e$ becomes $dm_e$ the rate of mass expulsion---or ``mass flow
rate''---is $dm_e/dt$ (also written $\mdot$). The product of the mass flow rate and the exhaust velocity
produces a component of force $-\vec{c}(dm_e/dt)$ acting in the $(+y)$ direction. The exhaust gases leave
the rocket at the exhaust pressure $P_e$ when they cross the plane at the aft end of the rocket motor
nozzle, called the \emph{nozzle exit plane}. The difference between $P_e$ and the pressure of the atmosphere
outside the exhaust stream (also called the \emph{ambient} pressure) $P_a$ acts over the area of the nozzle
exit plane $\vec{A}_e$ to produce another component of force $\overrightarrow{(P_e - P_a)A_e}$ which also
acts in the $(+y)$ direction when $P_e$ is greater than $P_a$, but which acts in the $(-y)$ direction when
$P_e$ is less than $P_a$. The sum of the two force components is the rocket's \emph{thrust} $\vec{F}$:
\begin{gather*}
\vec{F} = -\vec{c}(dm_e/dt) + \overrightarrow{(P_e - P_a)A_e}
\end{gather*}
\emph{NOTE}: In order to make the illustrations easier to see, the exhaust stream has been drawn as if the
nozzle exit diameter were equal to the body tube diameter. The nozzle exit of most model rocket engines is
actually much smaller than the body tube diameter.

\bigskip
\noindent\begin{minipage}{\linewidth}
  \edcap{The 1973 Figure 2, which this figure replaces:}

  \begin{center}
    \includegraphics{figures/v2/ch1/fig02.pdf}
  \end{center}

  Figure 2: Origin of rocket thrust. In a time interval $\Delta t$ the rocket expels the quantity of mass
  $\Delta m_e$ at the exhaust velocity $\vec{c}$ in the $(-y)$ direction. ``In the limit'' as $\Delta t$
  becomes $dt$ and $\Delta m_e$ becomes $dm_e$ the rate of mass expulsion---or ``mass flow rate''---is
  $dm_e/dt$ (also written $\mdot$). The thrust $\vec{F}$ is given by the equation
  $\vec{F} = -\vec{c}(dm_e/dt)$ and thus acts in the $(+y)$ direction.
\end{minipage}

\chapter{Correction to Original Page 40 of Chapter 1}\label{supp:ch1-page40}

According to the method of analysis developed by James Barrowman, which is discussed in detail in
Chapter~\ref{ch2}, a rocket flying at a nonzero angle of attack experiences a \emph{normal force} (i.e., one
perpendicular to the model's centerline) whose magnitude is given by
\begin{equation*}\tag{20}
N = (1/2)\rho C_N A_r v^{2}
\end{equation*}
where $C_N$ is the \emph{normal force coefficient} whose computation by the Barrowman method is described in
Chapter~\ref{ch2}. For angles of attack that are not too large $C_N$ is very nearly directly proportional to
$\alpha$. The variation of $C_N$ with $\alpha$ can then be expressed as
\begin{equation*}\tag{21}
C_N = \CNa \cdot \alpha
\end{equation*}
where $\CNa$ is the slope of the curve of the normal force coefficient versus angle of attack, often referred
to in an abbreviated fashion as the \emph{normal force curve slope}. Then
\begin{equation*}\tag{22}
N = (1/2)\rho \CNa A_r v^{2}\alpha
\end{equation*}
The component of force $N\sin\alpha$ is an additional contribution to the \emph{drag}, while the component of
force $N\cos\alpha$ is a \emph{side force} tending to deflect the rocket's path by causing an acceleration
component normal to the current direction of flight. The resolution of the normal force is shown in
Figure~\ref{ch1:fig:7}.

\bigskip
\begin{center}
[All subsequent equation numbers must be increased by 2]
\end{center}
